\documentclass[12pt,a4paper]{article}

% ------------------------------------------------------------
% Encoding, language, typography
% ------------------------------------------------------------
\usepackage[utf8]{inputenc}
\usepackage[T1]{fontenc}
\usepackage[english]{babel}
\usepackage{lmodern}
\usepackage{microtype}

% ------------------------------------------------------------
% Page layout
% ------------------------------------------------------------
\usepackage[a4paper,margin=2.5cm]{geometry}
\usepackage{setspace}
\onehalfspacing

% ------------------------------------------------------------
% Mathematics
% ------------------------------------------------------------
\usepackage{amsmath}
\usepackage{amssymb}
\usepackage{amsfonts}
\usepackage{mathtools}
\usepackage{bm}

% ------------------------------------------------------------
% Tables, lists, links
% ------------------------------------------------------------
\usepackage{booktabs}
\usepackage{enumitem}
\usepackage{xcolor}
\usepackage[hidelinks]{hyperref}

% ------------------------------------------------------------
% Citations
% ------------------------------------------------------------
\usepackage[round,authoryear]{natbib}

% ------------------------------------------------------------
% Theorem-style environments, optional
% ------------------------------------------------------------
\usepackage{amsthm}
\newtheorem{assumption}{Assumption}
\newtheorem{definition}{Definition}
\newtheorem{proposition}{Proposition}
\newtheorem{remark}{Remark}

% ------------------------------------------------------------
% Useful math commands
% ------------------------------------------------------------
\newcommand{\E}{\mathbb{E}}
\newcommand{\En}{\mathbb{E}_N}
\newcommand{\R}{\mathbb{R}}
\newcommand{\ind}{\mathbbm{1}}

\title{Partially Linear Instrumental Variables and Double Machine Learning}
\author{}
\date{}

\begin{document}

\maketitle

\section{Double Machine Learning for Instrumental Variables}

\paragraph{Assumptions for the partially linear IV model.}
The partially linear instrumental variables model relies on three distinct types of assumptions:
(i) structural/modeling assumptions, (ii) identification assumptions, and (iii) estimation
assumptions required for Double Machine Learning.

\medskip
\noindent
\textbf{Structural/modeling assumption.}
The outcome equation is partially linear:
\begin{align}
    Y_i = D_i\theta_0 + g_0(X_i) + U_i.
    \label{eq:pliv_structural_assumption}
\end{align}
The treatment enters the outcome equation linearly with constant coefficient \(\theta_0\),
while the covariate component \(g_0(X_i)\) is left unrestricted. Thus, the model relaxes
linearity in the controls but does not relax effect homogeneity.

\medskip
\noindent
\textbf{Instrument validity conditional on covariates.}
The instrument is valid conditional on \(X_i\):
\begin{align}
    \mathbb{E}[U_i \mid X_i,Z_i] = 0.
    \label{eq:pliv_conditional_exogeneity}
\end{align}
This replaces the non-IV partially linear regression condition
\(\mathbb{E}[U_i \mid X_i,D_i]=0\). Hence, \(D_i\) may be endogenous, but the instrument
\(Z_i\) must be orthogonal to the structural error after conditioning on \(X_i\).

\medskip
\noindent
\textbf{Instrument residualization.}
The instrument can be decomposed into a covariate-predictable part and a residual component:
\begin{align}
    Z_i = m_0(X_i) + V_i,
    \qquad
    \mathbb{E}[V_i \mid X_i] = 0.
    \label{eq:pliv_instrument_residualization}
\end{align}
Equivalently,
\begin{align}
    V_i = Z_i - m_0(X_i),
    \qquad
    m_0(X_i)=\mathbb{E}[Z_i\mid X_i].
\end{align}
The residual \(V_i\) is the part of the instrument that remains after partialling out the
covariates.

\medskip
\noindent
\textbf{Residualized first-stage relevance.}
The residualized instrument must predict the treatment:
\begin{align}
    \mathbb{E}[\widetilde{Z}_i\widetilde{D}_i]\neq 0,
    \label{eq:pliv_residualized_first_stage}
\end{align}
where
\begin{align}
    \widetilde{Z}_i &= Z_i-\mathbb{E}[Z_i\mid X_i], \\
    \widetilde{D}_i &= D_i-\mathbb{E}[D_i\mid X_i].
\end{align}
This is the analogue of the usual IV relevance condition after removing the variation explained
by \(X_i\). In the notation of \citet{chernozhukov2018double}, this condition corresponds to
\(|\mathbb{E}[D_iV_i]|\geq c>0\).

\medskip
\noindent
\textbf{Nuisance functions.}
For the Robinson-style PLIV score, define the nuisance functions
\begin{align}
    \ell_0(X_i) &= \mathbb{E}[Y_i\mid X_i], \\
    r_0(X_i) &= \mathbb{E}[D_i\mid X_i], \\
    m_0(X_i) &= \mathbb{E}[Z_i\mid X_i].
\end{align}
The corresponding residualized variables are
\begin{align}
    \widetilde{Y}_i &= Y_i-\ell_0(X_i), \\
    \widetilde{D}_i &= D_i-r_0(X_i), \\
    \widetilde{Z}_i &= Z_i-m_0(X_i).
\end{align}

\medskip
\noindent
\textbf{Orthogonal moment condition.}
The identifying moment condition is
\begin{align}
    \mathbb{E}
    \left[
        \left(
            \widetilde{Y}_i-\theta_0\widetilde{D}_i
        \right)
        \widetilde{Z}_i
    \right]
    =0.
    \label{eq:pliv_orthogonal_moment}
\end{align}
Equivalently, the Neyman-orthogonal score is
\begin{align}
    \psi(W_i;\theta,\eta)
    =
    \left(
        Y_i-\ell(X_i)-\theta(D_i-r(X_i))
    \right)
    \left(
        Z_i-m(X_i)
    \right),
    \qquad
    \eta=(\ell,r,m).
    \label{eq:pliv_score_assumption_block}
\end{align}
Neyman orthogonality means that small first-stage errors in the nuisance functions affect the
final estimator only in second order.

\medskip
\noindent
\textbf{DML estimation assumptions.}
The nuisance functions \(\ell_0\), \(r_0\), and \(m_0\) are estimated by machine learning methods
using cross-fitting. Cross-fitting requires that the prediction for observation \(i\) is obtained from
models not trained on observation \(i\). In addition, the nuisance estimators must be sufficiently
accurate. A sufficient product-rate condition for the Robinson-style PLIV score is
\begin{align}
    \|\hat m-m_0\|_{P,2}
    \left(
        \|\hat r-r_0\|_{P,2}
        +
        \|\hat\ell-\ell_0\|_{P,2}
    \right)
    =
    o(N^{-1/2}).
    \label{eq:pliv_product_rate_assumption_block}
\end{align}
Together with standard moment and boundedness conditions, this yields root-\(N\) consistent
and asymptotically normal inference for \(\theta_0\).

\subsection{Partially Linear Instrumental Variables and Double Machine Learning}
\label{subsec:pliv_dml}

The double/debiased machine learning framework of \citet{chernozhukov2018double}
provides a way to combine flexible nuisance-function estimation with valid inference
on a low-dimensional target parameter. A useful starting point is the partially linear
regression model of \citet{robinson1988root},
\begin{align}
    Y_i = D_i \theta_0 + g_0(X_i) + U_i,
    \qquad
    \E[U_i \mid X_i,D_i] = 0,
    \label{eq:plr_model}
\end{align}
where \(Y_i\) denotes the outcome, \(D_i\) the treatment variable, \(X_i\) a vector of
covariates, and \(g_0(\cdot)\) an unknown nuisance function. The parameter of interest
is the scalar coefficient \(\theta_0\). The model relaxes the linear specification of the
covariates by allowing \(g_0(X_i)\) to be an unrestricted function, but it still imposes a
constant treatment effect \(\theta_0\).

In the applications considered in this thesis, however, \(D_i\) is not assumed to be
conditionally exogenous. Instead, identification relies on an instrumental variable
\(Z_i\). Following \citet[Section~4.2]{chernozhukov2018double}, the partially linear
instrumental variables model extends \eqref{eq:plr_model} by replacing conditional
exogeneity of \(D_i\) with conditional validity of the instrument:
\begin{align}
    Y_i &= D_i \theta_0 + g_0(X_i) + U_i,
    \qquad
    \E[U_i \mid X_i,Z_i] = 0,
    \label{eq:pliv_structural} \\
    Z_i &= m_0(X_i) + V_i,
    \qquad
    \E[V_i \mid X_i] = 0.
    \label{eq:pliv_firststage_z}
\end{align}
Here, \(Z_i\) is the instrumental variable and \(V_i = Z_i - m_0(X_i)\) denotes the part
of the instrument that is not predictable from the covariates. Thus, the model allows
\(D_i\) to be endogenous, but requires the residualized instrument to be orthogonal to
the structural error.


The moment condition implied by \eqref{eq:pliv_structural}--\eqref{eq:pliv_firststage_z}
can be written as
\begin{align}
    \E
    \left[
        \left(Y_i - D_i \theta_0 - g_0(X_i)\right)
        \left(Z_i - m_0(X_i)\right)
    \right]
    = 0.
    \label{eq:pliv_moment_g}
\end{align}
Equivalently, using the notation \(W_i=(Y_i,D_i,X_i,Z_i)\), this corresponds to the
Neyman-orthogonal score
\begin{align}
    \psi(W_i;\theta,\eta)
    =
    \left(Y_i - D_i\theta - g(X_i)\right)
    \left(Z_i - m(X_i)\right),
    \qquad
    \eta=(g,m).
    \label{eq:pliv_score_g}
\end{align}

For implementation, \citet{chernozhukov2018double} also propose a Robinson-style
partialling-out score. Define the nuisance functions
\begin{align}
    \ell_0(X_i) &= \E[Y_i \mid X_i], \label{eq:nuisance_l} \\
    r_0(X_i) &= \E[D_i \mid X_i], \label{eq:nuisance_r} \\
    m_0(X_i) &= \E[Z_i \mid X_i]. \label{eq:nuisance_m}
\end{align}
Then the residualized variables are
\begin{align}
    \widetilde{Y}_i &= Y_i - \ell_0(X_i), \label{eq:y_tilde_pliv} \\
    \widetilde{D}_i &= D_i - r_0(X_i), \label{eq:d_tilde_pliv} \\
    \widetilde{Z}_i &= Z_i - m_0(X_i). \label{eq:z_tilde_pliv}
\end{align}
The corresponding orthogonal score is
\begin{align}
    \psi(W_i;\theta,\eta)
    =
    \left(
        Y_i - \ell(X_i)
        -
        \theta\left(D_i-r(X_i)\right)
    \right)
    \left(
        Z_i-m(X_i)
    \right),
    \qquad
    \eta=(\ell,r,m).
    \label{eq:pliv_robinson_score}
\end{align}
At the true nuisance functions, this can be written compactly as
\begin{align}
    \psi(W_i;\theta_0,\eta_0)
    =
    \left(
        \widetilde{Y}_i
        -
        \theta_0 \widetilde{D}_i
    \right)
    \widetilde{Z}_i.
    \label{eq:pliv_residual_score}
\end{align}
The identifying moment condition is therefore
\begin{align}
    \E
    \left[
        \left(
            \widetilde{Y}_i
            -
            \theta_0 \widetilde{D}_i
        \right)
        \widetilde{Z}_i
    \right]
    = 0.
    \label{eq:pliv_residual_moment}
\end{align}

Solving \eqref{eq:pliv_residual_moment} gives the population coefficient
\begin{align}
    \theta_0
    =
    \frac{
        \E[\widetilde{Z}_i \widetilde{Y}_i]
    }{
        \E[\widetilde{Z}_i \widetilde{D}_i]
    },
    \label{eq:pliv_population_ratio}
\end{align}
provided that the residualized first stage satisfies
\begin{align}
    \E[\widetilde{Z}_i \widetilde{D}_i] \neq 0.
    \label{eq:pliv_relevance}
\end{align}
This condition is the analogue of instrument relevance after partialling out the covariates.

In practice, the nuisance functions \(\ell_0\), \(r_0\), and \(m_0\) are unknown and are
estimated flexibly by machine learning methods. Let \(\hat{\ell}\), \(\hat r\), and
\(\hat m\) denote cross-fitted estimates of the nuisance functions. The DML estimator based
on the Robinson-style score is then
\begin{align}
    \hat{\theta}_{\mathrm{PLIV}}
    =
    \frac{
        \sum_{i=1}^N
        \left(Z_i-\hat m(X_i)\right)
        \left(Y_i-\hat{\ell}(X_i)\right)
    }{
        \sum_{i=1}^N
        \left(Z_i-\hat m(X_i)\right)
        \left(D_i-\hat r(X_i)\right)
    }.
    \label{eq:pliv_dml_estimator}
\end{align}
Thus, the partially linear IV estimator can be interpreted as an instrumental variables
estimator applied to residualized outcome, treatment, and instrument variables.

The theoretical validity of this estimator rests on three types of conditions. First, the
partially linear IV model \eqref{eq:pliv_structural}--\eqref{eq:pliv_firststage_z} must hold.
Second, the residualized first stage in \eqref{eq:pliv_relevance} must be nonzero. Third,
the nuisance estimators must be sufficiently accurate. In the DML framework, this requirement
is expressed through product-rate conditions on the first-stage errors. For the Robinson-style
score, a sufficient condition is of the form
\begin{align}
    \left\|
        \hat m - m_0
    \right\|_{P,2}
    \left(
        \left\|
            \hat r-r_0
        \right\|_{P,2}
        +
        \left\|
            \hat{\ell}-\ell_0
        \right\|_{P,2}
    \right)
    =
    o(N^{-1/2}),
    \label{eq:pliv_product_rate}
\end{align}
up to the precise regularity conditions stated in \citet[Assumption~4.2]{chernozhukov2018double}.
Together with cross-fitting and Neyman orthogonality, these conditions imply root-\(N\)
consistent and asymptotically normal inference for \(\theta_0\).


\subsection{Distinguishing PLIV from LATE Estimators}
\label{subsec:pliv_vs_late}

The estimators considered in this thesis are all related to instrumental variables, but they do
not all target the same causal parameter under the same assumptions. This distinction is
important when comparing kappa-based estimators, Wald-AIPW, and PLIV.

The main target parameter in the kappa-weighting framework is the local average treatment
effect (LATE). Under the standard IV assumptions---conditional instrument independence,
exclusion, instrument relevance, overlap, and monotonicity---the LATE is defined as
\begin{align}
    \tau_{\mathrm{LATE}}
    =
    \mathbb{E}
    \left[
        Y_i(1)-Y_i(0)
        \mid
        D_i(1)>D_i(0)
    \right].
    \label{eq:late_definition_comparison}
\end{align}
Thus, the LATE is the average treatment effect for compliers, i.e. for those individuals whose
treatment status is changed by the instrument. Kappa-based estimators and the DML
Wald-AIPW estimator are both naturally interpreted as estimators of this complier-average
causal effect. They differ in how they adjust for covariates and in the outcome weights they
imply, but they are designed to target the same LATE parameter.

The partially linear IV estimator, by contrast, is based on a different modeling structure. In
the PLIV model, the outcome equation is written as
\begin{align}
    Y_i
    =
    D_i\theta_0
    +
    g_0(X_i)
    +
    U_i,
    \qquad
    \mathbb{E}[U_i\mid X_i,Z_i]=0.
    \label{eq:pliv_structural_comparison}
\end{align}
Here, \(\theta_0\) is the constant coefficient on the endogenous treatment \(D_i\), while
\(g_0(X_i)\) is an unknown and potentially flexible function of the covariates. Hence, PLIV
relaxes the functional form of the covariate adjustment, but it does not relax effect
homogeneity. The model imposes that the treatment enters the structural outcome equation
through a single scalar coefficient \(\theta_0\).

Using the Robinson-style partialling-out representation, define

which yields the population ratio
\begin{align}
    \theta_0
    =
    \frac{
        \mathbb{E}[\widetilde{Z}_i\widetilde{Y}_i]
    }{
        \mathbb{E}[\widetilde{Z}_i\widetilde{D}_i]
    }.
    \label{eq:pliv_ratio_comparison}
\end{align}
This expression has a Wald-type form, since it divides a residualized reduced-form relationship
by a residualized first-stage relationship. However, it should not be confused with the LATE
Wald estimand. The PLIV ratio estimates the constant structural IV coefficient in
\eqref{eq:pliv_structural_comparison}, whereas the LATE estimand in
\eqref{eq:late_definition_comparison} averages heterogeneous treatment effects over the
complier subpopulation.

The two parameters coincide only under additional effect-homogeneity restrictions. If
\begin{align}
    Y_i(1)-Y_i(0)=\theta_0
    \qquad
    \text{for all } i,
    \label{eq:homogeneous_effects_comparison}
\end{align}
then the average treatment effect, the local average treatment effect, and the PLIV coefficient
all reduce to the same scalar parameter. Without such a homogeneity restriction, however, the
PLIV coefficient should be interpreted as a structural IV benchmark rather than as a direct
LATE estimator.

Therefore, the empirical comparison in this thesis distinguishes between two groups of
estimators. PLIV is reported separately as an additional DML-IV benchmark. It is useful
because it shares the residualized pseudo-IV structure and admits an outcome-weight
representation, but it does not generally target the same parameter as the kappa LATE
estimators under treatment-effect heterogeneity.
% ------------------------------------------------------------
% Bibliography placeholders
% ------------------------------------------------------------
\bibliographystyle{apalike}
\bibliography{references}

\end{document}